%%
%% Datei: paper73_erdos_ko_rado_de.tex
%% Autor: Michael Fuhrmann
%% Beschreibung: Der Erdős-Ko-Rado-Satz: Verallgemeinerungen und offene Probleme
%%               (Deutsche Version)
%% Erstellt: 2026-03-12
%% lastModified: 2026-03-12
%% Build: 164
%%
\documentclass[12pt,a4paper]{amsart}

\usepackage{amsmath,amssymb,amsthm,mathtools,enumitem,booktabs,array,hyperref}
\usepackage[utf8]{inputenc}
\usepackage[T1]{fontenc}
\usepackage[ngerman]{babel}
\usepackage{geometry}
\geometry{margin=2.5cm}

%% Theorem-Umgebungen (Englisch, Basis)
\newtheorem{theorem}{Theorem}[section]
\newtheorem{lemma}[theorem]{Lemma}
\newtheorem{corollary}[theorem]{Corollary}
\newtheorem{proposition}[theorem]{Proposition}
\theoremstyle{definition}
\newtheorem{definition}[theorem]{Definition}
\newtheorem{example}[theorem]{Example}
\theoremstyle{remark}
\newtheorem{remark}[theorem]{Remark}
\newtheorem{conjecture}[theorem]{Conjecture}

%% Deutsche Theorem-Umgebungen
\newtheorem{satz}[theorem]{Satz}
\newtheorem{vermutung}[theorem]{Vermutung}
\newtheorem{korollar}[theorem]{Korollar}
\newtheorem{lemma_de}[theorem]{Lemma}
\newtheorem{bemerkung}[theorem]{Bemerkung}
\newtheorem{definition_de}[theorem]{Definition}

%% Kürzel
\newcommand{\F}{\mathbb{F}}
\newcommand{\Z}{\mathbb{Z}}
\newcommand{\N}{\mathbb{N}}

\begin{document}
%% Abstract VOR \maketitle
\begin{abstract}
Der Erdős--Ko--Rado-Satz (EKR) von 1961 ist das grundlegende Ergebnis der
extremalen Mengenlehre: Die größte schneidende $k$-uniforme Familie über einer
$n$-elementigen Grundmenge ist ein \emph{Stern} (alle Mengen enthalten ein festes
Element), sofern $n \geq 2k$. Wir geben eine Gesamtdarstellung des Originalsatzes
mit Katonas elegantem Kreisbeweis, der Hilton--Milner-Charakterisierung der
zweitgrößten Familie, Frankls $t$-schneidender Verallgemeinerung, dem vollständigen
Schnittsatz von Ahlswede--Khachatrian und den $q$-Analoga über Vektorräumen.
Kreuzschneidende Familien, Zeichenmengen und Permutationsgruppen werden behandelt.
Frankls Vereinigungs-geschlossene Vermutung und mehrere Extremalprobleme für
nicht-uniforme Hypergraphen bleiben offen.
\end{abstract}

\title{Der Erd\H{o}s--Ko--Rado-Satz:\\
Verallgemeinerungen und offene Probleme}

\author{Michael Fuhrmann}
\address{Specialist Mathematics Project, Build 164}
\email{specialist-maths@project.local}
\date{2026-03-12}

\maketitle

\tableofcontents

%% ============================================================
\section{Einleitung}
%% ============================================================

Sei $[n] = \{1, 2, \ldots, n\}$ und $\binom{[n]}{k}$ die Menge aller
$k$-elementigen Teilmengen von $[n]$. Eine Familie
$\mathcal{F} \subseteq \binom{[n]}{k}$ heißt \emph{schneidend} (intersecting), wenn
$A \cap B \neq \emptyset$ für alle $A, B \in \mathcal{F}$.

Die einfachste schneidende Familie ist der \emph{Stern}
$\mathcal{S}_i = \{A \in \binom{[n]}{k} : i \in A\}$
mit Größe $\binom{n-1}{k-1}$. Die grundlegende Frage lautet: Ist dies die größte
mögliche schneidende Familie?

\begin{satz}[Erdős--Ko--Rado \cite{ErdosKoRado1961}]\label{satz:ekr}
Sei $n \geq 2k$ und $\mathcal{F} \subseteq \binom{[n]}{k}$ schneidend. Dann gilt:
\[
  |\mathcal{F}| \leq \binom{n-1}{k-1}.
\]
Gleichheit gilt (für $n > 2k$) genau dann, wenn $\mathcal{F}$ ein Stern ist.
\end{satz}

Dieser Satz hat eine enorme Literatur erzeugt. Seine Verallgemeinerungen berühren
Kombinatorik, Algebra, Codierungstheorie und diskrete Geometrie.

%% ============================================================
\section{Katonas Kreisbeweis}\label{sec:katona}
%% ============================================================

Der eleganteste Beweis des EKR-Satzes stammt von Katona \cite{Katona1972}.

\begin{proof}[Katonas Beweis von Satz~\ref{satz:ekr}]
Betrachte alle $n!$ zyklischen Permutationen von $[n]$ auf einem Kreis angeordnet.
Für jede $k$-elementige Menge $A \subseteq [n]$ erscheint sie als aufeinanderfolgender
Bogen in genau $k!(n-k)!$ dieser zyklischen Anordnungen.

Für eine feste zyklische Anordnung $\sigma$ heißt $A \in \mathcal{F}$ ein
\emph{guter Bogen}, wenn er als aufeinanderfolgender Bogen in $\sigma$ erscheint.
Da $\mathcal{F}$ schneidend ist, müssen je zwei gute Bögen in $\sigma$ sich
überlappen; da sie alle Länge $k$ auf einem Kreis der Länge $n \geq 2k$ haben,
können höchstens $k$ der $n$ möglichen $k$-Bögen in $\sigma$ paarweise schneiden
(sie müssen sich alle um einen gemeinsamen Punkt "stapeln"). Also erscheinen
höchstens $k$ Elemente von $\mathcal{F}$ als Bögen in $\sigma$.

Zählen der Paare $(\sigma, A)$ mit $A$ als gutem Bogen von $\sigma$:
\[
  \sum_\sigma |\{\text{gute Bögen von } \sigma \text{ in } \mathcal{F}\}|
  = |\mathcal{F}| \cdot k!(n-k)!.
\]
Da jedes $\sigma$ höchstens $k$ beiträgt:
\[
  |\mathcal{F}| \cdot k!(n-k)! \leq k \cdot n!,
  \quad\Longrightarrow\quad
  |\mathcal{F}| \leq \frac{k \cdot n!}{k!(n-k)!} = \binom{n-1}{k-1}.
\qedhere\]
\end{proof}

%% ============================================================
\section{Der Hilton--Milner-Satz}\label{sec:hilton_milner}
%% ============================================================

\begin{satz}[Hilton--Milner \cite{HiltonMilner1967}]\label{satz:hilton_milner}
Sei $n > 2k$ und $\mathcal{F} \subseteq \binom{[n]}{k}$ schneidend mit
$\bigcap_{A \in \mathcal{F}} A = \emptyset$ (kein gemeinsames Element). Dann gilt:
\[
  |\mathcal{F}| \leq \binom{n-1}{k-1} - \binom{n-k-1}{k-1} + 1.
\]
\end{satz}

\begin{proof}[Beweisskizze]
Da $\mathcal{F}$ schneidend, aber ohne gemeinsames Element ist, gibt es
$A, B \in \mathcal{F}$ mit $A \cap B = \{x\}$ für ein $x$ (das "Zeugenpaar").
Jedes weitere $F \in \mathcal{F}$ muss sowohl $A$ als auch $B$ schneiden, also
entweder $x \in F$ oder $F$ schneidet $A \setminus \{x\}$ und $B \setminus \{x\}$.
Sorgfältiges Abzählen dieser Fälle liefert die Schranke.
\end{proof}

\begin{bemerkung}
Die extremale Familie für Hilton--Milner ist
$\mathcal{HM}(n,k) = \{A \in \binom{[n]}{k} : 1 \in A,\, A \cap [k+1] \neq \emptyset\}
\cup \{[2, k+1]\}$.
\end{bemerkung}

%% ============================================================
\section{Frankls $t$-schneidender Satz}\label{sec:frankl}
%% ============================================================

\begin{definition_de}[$t$-schneidende Familie]
Eine Familie $\mathcal{F} \subseteq \binom{[n]}{k}$ heißt \emph{$t$-schneidend},
wenn $|A \cap B| \geq t$ für alle $A, B \in \mathcal{F}$.
\end{definition_de}

\begin{satz}[Frankl \cite{Frankl1977}]\label{satz:frankl_t}
Sei $n$ hinreichend groß (Frankl-Schranke) und $\mathcal{F} \subseteq \binom{[n]}{k}$
$t$-schneidend. Dann gilt:
\[
  |\mathcal{F}| \leq \binom{n-t}{k-t}.
\]
\end{satz}

Das extremale Beispiel ist der \emph{$t$-Stern}:
$\mathcal{F}^* = \{A \in \binom{[n]}{k} : [t] \subseteq A\}$ mit Größe
$\binom{n-t}{k-t}$.

%% ============================================================
\section{Der vollständige Satz von Ahlswede--Khachatrian}\label{sec:ak}
%% ============================================================

Für alle Werte von $n, k, t$ wurde die vollständige Antwort von Ahlswede und
Khachatrian gegeben.

\begin{definition_de}[$r$-Kern-Familie]
Für $0 \leq r \leq k-t$ definiere die \emph{$r$-Kern-Familie}
\[
  \mathcal{A}_r = \{A \in \binom{[n]}{k} : |A \cap [t+2r]| \geq t+r\}.
\]
\end{definition_de}

\begin{satz}[Ahlswede--Khachatrian \cite{AhlswedeKhachatrian1997}]\label{satz:ak}
Sei $1 \leq t \leq k \leq n$ und $\mathcal{F} \subseteq \binom{[n]}{k}$
$t$-schneidend. Dann gilt:
\[
  |\mathcal{F}| \leq \max_{r \geq 0} |\mathcal{A}_r|,
\]
und Gleichheit gilt genau dann, wenn $\mathcal{F}$ isomorph zu einem $\mathcal{A}_r$ ist.
\end{satz}

\begin{bemerkung}
Dies ist eine vollständige Charakterisierung. Das Maximum wird beim durch $n, k, t$
bestimmten Wert von $r$ angenommen. Für $n$ groß relativ zu $k$ liegt das Maximum
bei $r=0$ (Frankls Satz).
\end{bemerkung}

%% ============================================================
\section{Kreuzschneidende Familien}\label{sec:kreuz}
%% ============================================================

\begin{definition_de}[Kreuzschneidend]
Zwei Familien $\mathcal{A}, \mathcal{B} \subseteq \binom{[n]}{k}$ heißen
\emph{kreuzschneidend}, wenn $A \cap B \neq \emptyset$ für alle $A \in \mathcal{A}$,
$B \in \mathcal{B}$.
\end{definition_de}

\begin{satz}[Pyber \cite{Pyber1986}]\label{satz:pyber}
Sind $\mathcal{A}, \mathcal{B} \subseteq \binom{[n]}{k}$ kreuzschneidend mit
$n \geq 2k$, so gilt:
\[
  |\mathcal{A}| \cdot |\mathcal{B}| \leq \binom{n-1}{k-1}^2.
\]
\end{satz}

\begin{vermutung}[Kreuzschneidende Summenvermutung]\label{verm:kreuz_summe}
Sind $\mathcal{A}, \mathcal{B} \subseteq \binom{[n]}{k}$ kreuzschneidend mit
$n \geq 2k$, so gilt:
\[
  |\mathcal{A}| + |\mathcal{B}| \leq \binom{n}{k} - \binom{n-k}{k} + 1.
\]
\end{vermutung}

%% ============================================================
\section{$q$-Analog: Vektorräume über $\mathbb{F}_q$}\label{sec:q_analog}
%% ============================================================

\begin{definition_de}[$q$-Analog]
Sei $V = \F_q^n$ der $n$-dimensionale Vektorraum über $\F_q$. Der
\emph{Gauß'sche Binomialkoeffizient} $\binom{n}{k}_q$ zählt die $k$-dimensionalen
Unterräume von $V$:
\[
  \binom{n}{k}_q = \prod_{i=0}^{k-1} \frac{q^{n-i} - 1}{q^{i+1} - 1}.
\]
Eine Familie $\mathcal{F}$ von $k$-dimensionalen Unterräumen heißt
\emph{schneidend}, wenn $\dim(U \cap W) \geq 1$ für alle $U, W \in \mathcal{F}$.
\end{definition_de}

\begin{satz}[$q$-Analog EKR \cite{FranklWilson1986}]\label{satz:q_ekr}
Sei $n \geq 2k$ und $\mathcal{F}$ eine schneidende Familie von $k$-dimensionalen
Unterräumen von $\F_q^n$. Dann gilt:
\[
  |\mathcal{F}| \leq \binom{n-1}{k-1}_q.
\]
Gleichheit gilt genau dann, wenn $\mathcal{F}$ ein \emph{$q$-Stern} ist: alle
$k$-Räume, die einen festen 1-dimensionalen Unterraum enthalten.
\end{satz}

\begin{proof}[Beweisskizze]
Katonas Kreisbeweis lässt sich mit einem $q$-Analog des Permutationszählarguments
adaptieren. Alternativ nutzt die Frankl--Wilson-Methode die Darstellungstheorie
von $GL_n(\F_q)$.
\end{proof}

\begin{bemerkung}
Das $q$-Analog ist streng allgemeiner: $\binom{n-1}{k-1}_q \to \binom{n-1}{k-1}$
für $q \to 1$ (in geeignetem Sinne). Die Theorie der $q$-Analoga verbindet sich
mit Gebäuden, Codierungstheorie und zufälligen simplizialen Komplexen.
\end{bemerkung}

%% ============================================================
\section{Zeichenmengen und Permutationsgruppen}\label{sec:zeichen}
%% ============================================================

\begin{definition_de}[Vorzeichenmenge]
Eine \emph{Vorzeichen-$k$-Menge} auf $[n]$ ist ein Paar $(A, \epsilon)$, wobei
$A \in \binom{[n]}{k}$ und $\epsilon: A \to \{+, -\}$ eine Vorzeichenfunktion ist.
Zwei Vorzeichenmengen sind \emph{schneidend}, wenn sie ein vorzeichengleiches
Element teilen: es gibt $x \in A \cap B$ mit $\epsilon_A(x) = \epsilon_B(x)$.
\end{definition_de}

\begin{satz}[Vorzeichen-EKR \cite{DezaFranklFuredi1983}]\label{satz:zeichen_ekr}
Die größte schneidende Familie von Vorzeichen-$k$-Mengen auf $[n]$ (für $n \geq 2k$)
hat Größe $2^{k-1}\binom{n-1}{k-1}$.
\end{satz}

\begin{definition_de}[Schneidende Permutationsfamilie]
Eine Familie $\mathcal{F}$ von Permutationen von $[n]$ heißt \emph{schneidend},
wenn für je zwei $\sigma, \tau \in \mathcal{F}$ ein $i$ mit $\sigma(i) = \tau(i)$
existiert.
\end{definition_de}

\begin{satz}[Deza--Frankl \cite{DezaFranklPerm1983}]\label{satz:deza_frankl_perm}
Die größte schneidende Permutationsfamilie von $[n]$ hat Größe $(n-1)!$;
die extremalen Familien sind Nebenklassen von Punktstabilisatoren.
\end{satz}

%% ============================================================
\section{Nicht-uniforme Hypergraphen}\label{sec:nicht_uniform}
%% ============================================================

Für nicht-uniforme (gemischtgroße) schneidende Familien ist die Situation
komplizierter.

\begin{definition_de}[Schneidende Antikette]
Eine Familie $\mathcal{F} \subseteq 2^{[n]}$ heißt \emph{schneidende Antikette},
wenn $\mathcal{F}$ schneidend ist und keine Menge in $\mathcal{F}$ eine andere
enthält.
\end{definition_de}

\begin{satz}[Bollobás-Mengenpaare-Ungleichung \cite{Bollobas1965}]\label{satz:bollobas}
Seien $(A_i, B_i)_{i=1}^m$ Mengenpaare mit $A_i \cap B_j = \emptyset$ genau dann
wenn $i = j$, $|A_i| = a$, $|B_i| = b$. Dann gilt $m \leq \binom{a+b}{a}$.
\end{satz}

\begin{vermutung}[Nicht-uniformes EKR]\label{verm:nicht_uniform}
Für jede schneidende Antikette $\mathcal{F} \subseteq 2^{[n]}$ gibt es ein
Element $x \in [n]$, das in mindestens der Hälfte der Mengen in $\mathcal{F}$
enthalten ist.
\end{vermutung}

%% ============================================================
\section{Frankls Vereinigungs-geschlossene Vermutung}\label{sec:vereinigung}
%% ============================================================

\begin{vermutung}[Frankls Vereinigungs-geschlossene Vermutung \cite{Frankl1979}]
\label{verm:frankl_vg}
Sei $\mathcal{F}$ eine endliche vereinigungs-geschlossene Familie endlicher Mengen
(d.\,h.\ $A, B \in \mathcal{F}$ impliziert $A \cup B \in \mathcal{F}$) mit
$\mathcal{F} \neq \{\emptyset\}$. Dann gibt es ein Element $x$, das in mindestens
der Hälfte der Mengen in $\mathcal{F}$ enthalten ist.
\end{vermutung}

\begin{bemerkung}
Diese Vermutung ist seit 1979 bekannt und hat trotz enormer Anstrengungen keinen
Beweis erhalten. Wichtige Teilresultate:
\begin{enumerate}[label=(\roman*)]
  \item Knill: wahr für $|\mathcal{F}| \leq 46$ \cite{Knill1994}.
  \item Bošnjak--Marković: wahr, wenn die Grundmenge höchstens 12 Elemente hat
        \cite{BosnjacMarkovic2008}.
  \item Gilmer (2022): Es gibt stets ein $x$ in mindestens $(3-\sqrt{5})/2 \approx 38{,}2\%$
        der Mengen \cite{Gilmer2022} (Entropiemethode).
  \item Chase--Lovett (2024): Weitere Verfeinerung der Entropiemethode,
        Verbesserung der Konstantenschranke \cite{ChaseLovett2024}.
\end{enumerate}
\end{bemerkung}

%% ============================================================
\section{Stabilitätsergebnisse}\label{sec:stabilitaet}
%% ============================================================

Ein \emph{Stabilitätsergebnis} fragt: Wenn $|\mathcal{F}|$ nahe am Maximum liegt,
muss $\mathcal{F}$ (strukturell) nahe an der extremalen Familie sein?

\begin{satz}[Frankl Stabilität \cite{Frankl1987stability}]\label{satz:frankl_stab}
Sei $n > (2+\varepsilon)k$ und $\mathcal{F} \subseteq \binom{[n]}{k}$ schneidend mit
$|\mathcal{F}| \geq (1-\delta)\binom{n-1}{k-1}$. Dann gibt es $i \in [n]$ mit
\[
  |\mathcal{F} \setminus \mathcal{S}_i| \leq \delta' \binom{n-2}{k-2},
\]
wobei $\delta' \to 0$ für $\delta \to 0$.
\end{satz}

%% ============================================================
\section{Algebraische Methoden: Das Polynomverfahren}\label{sec:algebraisch}
%% ============================================================

\begin{satz}[Frankl--Wilson-Satz \cite{FranklWilson1981}]\label{satz:frankl_wilson}
Sei $L \subseteq \Z_{\geq 0}$ eine Menge von $s$ Resten modulo einer Primzahl $p$,
und $\mathcal{F} \subseteq 2^{[n]}$ eine Familie, sodass $|A| \not\equiv 0 \pmod{p}$
für alle $A \in \mathcal{F}$ und $|A \cap B| \equiv \lambda_i \pmod{p}$ für ein
$\lambda_i \in L$ gilt, wenn $A \neq B$. Dann gilt $|\mathcal{F}| \leq \binom{n}{s}$.
\end{satz}

\begin{korollar}[Untere Schranke für chromatische Zahl]\label{kor:chrom}
Der Kneser-Graph $K(n, k)$ (Knoten sind $k$-Mengen, Kanten verbinden disjunkte
Mengen) erfüllt $\chi(K(n, k)) \geq n - 2k + 2$. Dies bewies Lovász erstmals mit
topologischen Methoden \cite{Lovasz1978} und wurde später via den Borsuk--Ulam-Satz
neu bewiesen.
\end{korollar}

%% ============================================================
\section{Offene Probleme}\label{sec:offen}
%% ============================================================

\begin{vermutung}[Optimale kreuzschneidende Familien]\label{verm:kreuz_opt}
Für kreuzschneidende Familien $\mathcal{A}, \mathcal{B} \subseteq \binom{[n]}{k}$
mit $n \geq 2k$ gilt:
\[
  |\mathcal{A}| + |\mathcal{B}| \leq \binom{n}{k} - \binom{n-k}{k} + 1.
\]
\end{vermutung}

\begin{bemerkung}[Offenes Problem: Kurzer Beweis von Ahlswede--Khachatrian]
\label{bem:t_klein_n}
Der Satz von Ahlswede--Khachatrian (Satz~\ref{satz:ak}) bestimmt die maximale
Größe einer $t$-schneidenden Familie in $\binom{[n]}{k}$ vollständig für alle
$n, k, t$. Der Beweis ist jedoch lang und nicht-elementar; einen kurzen, konzeptuellen
Beweis zu finden, bleibt eine offene Herausforderung.
\end{bemerkung}

\begin{vermutung}[EKR für Permutationen mit $t \geq 2$]\label{verm:ekr_perm}
Die größte $t$-schneidende Permutationsfamilie von $[n]$ (mit Übereinstimmung in
$\geq t$ Stellen) hat Größe $(n-t)!$. Bekannt ist dies für $t=1$ (Deza--Frankl),
der allgemeine Fall ist offen.
\end{vermutung}

\begin{vermutung}[Frankl Vereinigungs-geschlossen, wiederholt]
\label{verm:frankl_vg2}
(Vgl.\ Vermutung~\ref{verm:frankl_vg}.) Dies gilt weithin als eines der
schönsten und zugänglichsten offenen Probleme der Kombinatorik.
\end{vermutung}

%% ============================================================
\section{Schluss}
%% ============================================================

Der Erdős--Ko--Rado-Satz ist eines der einflussreichsten Ergebnisse der
Kombinatorik. Sein Beweis durch Katonas Kreismethode ist ein Muster an Eleganz.
Die Verallgemeinerungen erstrecken sich auf $t$-schneidende Familien
(Frankl, Ahlswede--Khachatrian), kreuzschneidende Familien, $q$-Analoga über
Vektorräumen, Vorzeichenmengen und Permutationsgruppen. Die algebraische
Frankl--Wilson-Methode verbindet EKR mit Codierungstheorie und chromatischen
Zahlen.

Trotz mehr als 60 Jahren Forschung bleiben grundlegende Probleme offen: Frankls
Vereinigungs-geschlossene Vermutung, optimale kreuzschneidende Schranken und
EKR-Analoga für Permutationen mit $t \geq 2$. Die von Gilmer und Chase--Lovett
eingeführten Entropiemethoden eröffnen vielversprechende neue Richtungen.

%% Literaturverzeichnis
\begin{thebibliography}{99}

\bibitem{AhlswedeKhachatrian1997}
R.~Ahlswede, L.\,H.~Khachatrian,
\textit{The complete intersection theorem for systems of finite sets},
European J.\ Combin.\ \textbf{18} (1997), 125--136.

\bibitem{Bollobas1965}
B.~Bollobás,
\textit{On generalized graphs},
Acta Math.\ Acad.\ Sci.\ Hungar.\ \textbf{16} (1965), 447--452.

\bibitem{BosnjacMarkovic2008}
I.~Bošnjak, P.~Marković,
\textit{The 11-element case of Frankl's conjecture},
Electron.\ J.\ Combin.\ \textbf{15} (2008), R88.

\bibitem{ChaseLovett2024}
Z.~Chase, S.~Lovett,
\textit{Approximate union closed sets and union-closed conjecture},
Preprint arXiv:2201.01800 (2024, aktualisiert).

\bibitem{DezaFranklFuredi1983}
M.~Deza, P.~Frankl, Z.~Füredi,
\textit{The Erdős--Ko--Rado theorem for signed sets},
J.\ Combin.\ Inform.\ System Sci.\ \textbf{8} (1983), 87--92.

\bibitem{DezaFranklPerm1983}
M.~Deza, P.~Frankl,
\textit{Erdős--Ko--Rado theorem --- 22 years later},
SIAM J.\ Algebraic Discrete Methods \textbf{4} (1983), 419--431.

\bibitem{ErdosKoRado1961}
P.~Erd\H{o}s, C.~Ko, R.~Rado,
\textit{Intersection theorems for systems of finite sets},
Quart.\ J.\ Math.\ Oxford \textbf{12} (1961), 313--320.

\bibitem{Frankl1977}
P.~Frankl,
\textit{The Erdős--Ko--Rado theorem is true for $n = ckt$},
Colloq.\ Math.\ Soc.\ J.\ Bolyai \textbf{18} (1977), 365--375.

\bibitem{Frankl1979}
P.~Frankl,
\textit{An extremal problem for two families of sets},
European J.\ Combin.\ \textbf{3} (1979), 125--127.

\bibitem{Frankl1987stability}
P.~Frankl,
\textit{Stability of Erdős--Ko--Rado type results},
J.\ Combin.\ Theory Ser.\ A \textbf{46} (1987), 10--14.

\bibitem{FranklWilson1981}
P.~Frankl, R.\,M.~Wilson,
\textit{Intersection theorems with geometric consequences},
Combinatorica \textbf{1} (1981), 357--368.

\bibitem{FranklWilson1986}
P.~Frankl, R.\,M.~Wilson,
\textit{The Erdős--Ko--Rado theorem for vector spaces},
J.\ Combin.\ Theory Ser.\ A \textbf{43} (1986), 228--236.

\bibitem{Gilmer2022}
J.~Gilmer,
\textit{A constant lower bound for the union-closed conjecture},
Preprint arXiv:2211.09055 (2022).

\bibitem{HiltonMilner1967}
A.\,J.\,W.~Hilton, E.\,C.~Milner,
\textit{Some intersection theorems for systems of finite sets},
Quart.\ J.\ Math.\ Oxford \textbf{18} (1967), 369--384.

\bibitem{Katona1972}
G.\,O.\,H.~Katona,
\textit{A simple proof of the Erdős--Ko--Rado theorem},
J.\ Combin.\ Theory Ser.\ B \textbf{13} (1972), 183--184.

\bibitem{Knill1994}
E.~Knill,
\textit{Graph generated union-closed families of sets},
Preprint, Los Alamos National Laboratory, 1994.

\bibitem{Lovasz1978}
L.~Lov\'asz,
\textit{Kneser's conjecture, chromatic number, and homotopy},
J.\ Combin.\ Theory Ser.\ A \textbf{25} (1978), 319--324.

\bibitem{Pyber1986}
L.~Pyber,
\textit{A new generalization of the Erdős--Ko--Rado theorem},
J.\ Combin.\ Theory Ser.\ A \textbf{43} (1986), 85--90.

\bibitem{GodsilMeagherBook}
C.~Godsil, K.~Meagher,
\textit{Erd\H{o}s--Ko--Rado Theorems: Algebraic Approaches},
Cambridge University Press, Cambridge, 2016.

\end{thebibliography}

\end{document}
